\section{Experimental Setup}\label{sec:study}

To empirically evaluate the impact of a single \code{@Memoize} annotation on the runtime behaviour of existing code, we built a Jetpack Benchmark~\cite{jetpackbenchmark} harness around \numAlgorithms compute-intensive, overlapping-subproblem algorithms. Each algorithm is present in a baseline (un-memoized) variant and in three memoized variants, one per eviction policy described in Section~\ref{sec:policies}. Three benchmark suites operationalise the three research questions stated in Section~\ref{sec:intro}.

\subsection{Algorithms Under Test}\label{sec:algorithms}

The \numAlgorithms algorithms comprising the benchmark are canonical textbook dynamic-programming and recursion exercises, selected against three criteria. First, the set exercises a range of argument-tuple shapes (single integer, integer pair, and offsets into instance-held state) to ensure that the cache key construction path is exercised across shapes. Second, each algorithm is exponential in its naive recursive formulation and polynomial under memoization, so that the measured speedup is asymptotic rather than constant. Third, all arguments are primitive integers, which keeps the \code{Arrays.deepHashCode}-based cache-key construction path inexpensive and prevents key hashing from dominating the measured cost. Table~\ref{tab:algorithms} summarises the set.

\begin{table}[t]
\caption{The ten benchmark algorithms, with the asymptotic complexity reduction delivered by a single \code{@Memoize} annotation.}
\label{tab:algorithms}
\small
\centering
\begin{tabularx}{\columnwidth}{@{}l l l@{}}
\toprule
\textbf{Algorithm} & \textbf{Naive} & \textbf{Memoized} \\
\midrule
Fibonacci                 & $O(\varphi^n)$          & $O(n)$ \\
Ackermann                 & $\gg$ exponential       & O(reachable lattice) \\
Catalan                   & $O(C_n)$                & $O(n^2)$ \\
Binomial (Pascal)         & $O(2^n)$                & $O(n k)$ \\
Integer partition         & exponential             & $O(n k)$ \\
Edit distance             & $O(3^{n+m})$            & $O(n m)$ \\
Longest common subseq.    & $O(2^{n+m})$            & $O(n m)$ \\
Matrix chain              & $O(2^n)$                & $O(n^3)$ \\
0/1 knapsack              & $O(2^n)$                & $O(n W)$ \\
Coin change (ways)        & exponential             & $O(n A)$ \\
\bottomrule
\end{tabularx}
\end{table}


\paragraph{Primitive-key design.} For algorithms whose naive recursive formulation accepts array arguments (matrix-chain, knapsack, coin change) or string arguments (edit distance, longest common subsequence), the immutable input data is held as an instance field of the enclosing class, and only the indices into that data are exposed as method parameters. This ensures that the \code{CacheKeyWrapper} constructed by \code{MemoDispatcher.buildKey} contains only boxed \code{Integer} values, so cache-key construction remains O(1) in the size of the instance data.

\paragraph{Inheritance-based memoized variants.} To avoid duplicating ten method bodies three times across the three memoized variants, each subclass (\code{AlgorithmsLru}, \code{AlgorithmsFifo}, \code{AlgorithmsLfu}) extends the baseline \code{Algorithms} class and overrides every method with a body of the form \code{return super.method(args)}, carrying an \code{@Memoize(eviction~=~\ldots)} annotation. Because \code{javac} emits recursive self-calls inside \code{Algorithms} as \code{invokevirtual}, at runtime the recursion dispatches back to the overriding subclass method---and therefore re-enters the memoized cache-check sequence---causing the cache to cascade through the entire recursion tree. No source duplication is required across variants.

\subsection{Benchmark Suites}\label{sec:suites}

Each of the three benchmark suites operationalises exactly one research question.

\paragraph{Suite 1: \code{AlgorithmBenchmark}.} This suite answers \RQ{1}. It contains twenty test methods organised in pairs, one pair per algorithm: \code{<alg>\_baseline} measures a fresh \code{Algorithms} instance, and \code{<alg>\_memoized\_lru} measures a fresh \code{AlgorithmsLru} instance. A fresh instance is constructed for every \code{keepRunning} iteration, so that the measurement includes the cold-cache warm-up and the subsequent steady-state behaviour in a single window. This design matches how the un-memoized baseline measures the full recursion from scratch, making the per-iteration ratio a direct measure of the algorithmic speedup delivered by memoization.

\paragraph{Suite 2: \code{PolicyBenchmark}.} This suite answers \RQ{2}. It contains nine test methods exercising the three policies on three workloads of increasing eviction pressure: \code{fibonacci(25)} (narrow dense key space; the cache never overflows), \code{partition(40, 40)} (moderate two-dimensional key space; polynomial; fits the cache), and \code{coinChange(0, 120)} on a seven-denomination set (wider two-dimensional space, deliberately sized so that the cache is forced to evict). The first two workloads isolate \emph{pure per-operation cost}: eviction never happens, so all three policies return the same result and the only variation across benchmark rows is the implementation tax of the policy. The third introduces eviction pressure together with a skewed access distribution, which is the setting in which \lfu's theoretical advantage over \lru is expected to emerge.

\paragraph{Suite 3: \code{HitRateBenchmark}.} This suite answers \RQ{3}. It contains eleven test methods driving \code{fibonacci} through four workloads of decreasing reuse: \textbf{HOT} (five unique $n$ values repeated forty times each, yielding an approximate hit rate of 0.87); \textbf{WARM} (fifteen unique values drawn uniformly 200 times, yielding a steady-state hit rate of approximately 0.92 but with a larger working set); \textbf{COLD} (a Zipfian-like distribution in which 20\% of keys receive 80\% of the traffic, yielding an approximate hit rate of 0.80 with a long tail); and \textbf{VERY\_COLD} (a uniform distribution across a wide range of one-shot queries, yielding a hit rate close to zero). The baseline-to-memoized ratio is computed per regime; the point at which the ratio crosses 1.0 identifies the hit rate below which \code{@Memoize} costs more than it saves---the regime in which \code{autoMonitor} should intervene.

\subsection{Hardware and Profiling Infrastructure}\label{sec:hardware}

The harness is designed for execution on a physical Android device through \code{connectedReleaseAndroidTest}. Jetpack Benchmark refuses by default to report results collected on emulators, because thermal noise, CPU clock scaling, and host scheduler interference make emulator numbers unsuitable for citation. The target device class for this study is the Google Pixel family (Pixel 8 and Pixel 10) running stock Android 15 or 16; both devices are widely used in recent Android performance literature~\cite{empirical_rua} and expose \odpm power rails that support per-benchmark energy measurement in a follow-up study.

Following the methodology of the reference harness~\cite{empirical_rua}, the microbenchmark module uses \code{testBuildType~=~release} with R8 enabled, so that the reported numbers reflect production conditions including method inlining, devirtualisation, and dead-code elimination. A ProGuard rule file is shipped with the harness to preserve the \code{\_\_memoCacheManager} and \code{\_\_memoDispatcher\_*} synthetic fields during release-mode shrinking, which prevents R8 from stripping the memoization instrumentation. Jetpack Benchmark's built-in warm-up and \aot compilation machinery is used unmodified: multiple warm-up iterations precede every measurement block, the CPU clock is stabilised where possible, and the test target is \aot-compiled to eliminate \jit-driven variance. Every computed result is consumed via a \code{volatile long} sink field, preventing R8's full mode from eliminating the measured calls through dead-code analysis.

\subsection{Metrics}\label{sec:metrics}

For every test method, Jetpack Benchmark reports two quantities:

\begin{itemize}[leftmargin=*,nosep]
    \item \textbf{Time (ns).} Median per-iteration execution time with outlier trimming. This is the primary metric; all speedup ratios in Section~\ref{sec:results} are computed from this value.
    \item \textbf{Allocations (count).} Number of heap objects allocated per iteration. This metric serves as a proxy for garbage-collection pressure and allows the memory overhead of each memoized variant to be reported without attaching a full allocation profiler.
\end{itemize}

Every memoized class is annotated with \code{recordStats~=~true}. At the end of each benchmark run, \code{MemoCacheManager.dumpReport()} is emitted through \code{MemoLogger} at level \code{INFO}, producing per-dispatcher hit counts, miss counts, mean compute durations, mean lookup durations, and estimated savings. These statistics serve as a cross-check on the Jetpack Benchmark wall-clock numbers: a significant discrepancy between the expected and observed speedup would suggest a measurement artefact rather than an algorithmic effect.
